% Chapter 1 — Problems (§1.8)
\section{Problems}
\label{sec:ch01.8}

\noindent\emph{The exercises below assume harmonic fields under $\exp(+\jj\omega t)$ in SI units unless stated otherwise. Cite only results already established in this chapter or in standard vector-calculus identities used implicitly in the derivations.}

\begin{enumerate}
\item \textbf{Conservation laws and measurable fluxes.}
Assume a Maxwell-valid phasor pair $(\E,\H)$ on a control volume with the source and boundary data specified in Section~\ref{sec:ch01.1}.
\begin{enumerate}[label=(\alph*)]
\item Starting from $\mathbf{S}_\mathrm{c}=\tfrac12\E\times\H^*$, expand Eq.~\eqref{eq:ch1.poynting-real-imag-split} and state the physical meaning of $\Re(\mathbf{S}_\mathrm{c})$ and $\Im(\mathbf{S}_\mathrm{c})$ on a closed surface \cite{poynting1884,stratton1941}.
\item Using Eqs.~\eqref{eq:ch1.lorentz-force} and~\eqref{eq:ch1.maxwell-source-force}, show how contraction with Maxwell's equations yields Eq.~\eqref{eq:ch1.momentum-balance-instantaneous}. Identify when the volume force density $\mathbf{f}_{\mathrm{matter}}$ equals the impressed source force (Route~A) and when net body force is obtained from a source-free stress integral on $\mathcal{V}_{\mathrm{vac}}$ (Route~B).
\item Multiply the instantaneous angular-momentum balance by $\mathbf{r}\times$ (symmetric stress) to obtain Eq.~\eqref{eq:ch1.angular-momentum-balance}. Define $\boldsymbol{\mathcal{T}}_J(V)=\oint\Mavg\cdot\hat{\mathbf{n}}\,dS$ and $\boldsymbol{\Phi}_{J,\mathrm{out}}(V)=-\boldsymbol{\mathcal{T}}_J(V)$. \emph{Distinguish} Route~A source-support volume $V_{\mathrm{A}}$ from a material body volume $V_m$ and from the Route~B vacuum shell $\mathcal{V}_{\mathrm{vac}}$. For axis $\hat{\mathbf{u}}$, use the projected torque $\tau_{\hat{\mathbf{u}}}(V_{\mathrm{A}})$ of Eq.~\eqref{eq:ch1.torque-observable}---not $\tau_{\hat{\mathbf{u}}}(V_m)$ unless $V_m$ is explicitly declared to coincide with the Route~A enclosure. At harmonic steady state on a source-free enclosing surface, relate $\tau_{\hat{\mathbf{u}}}(V_{\mathrm{A}})$ to the outward stress-moment flux through $\hat{\mathbf{u}}\cdot\boldsymbol{\Phi}_{J,\mathrm{out}}$ when storage on $V_{\mathrm{A}}$ is stationary \cite{beth1936,stratton1941}.
\item A uniform plane wave of Eq.~\eqref{eq:ch1.plane-wave-transverse} is normally incident from $z<0$ on a perfect absorber at $z=0$ with outward vacuum normal $\hat{\mathbf{n}}=+\hat{\mathbf{z}}$. Verify $p_{\mathrm{abs}}=|\langle\mathbf{S}\rangle|/c$.
\item For a perfect reflector at normal incidence, verify $p_{\mathrm{refl}}=2|\langle\mathbf{S}\rangle|/c$ and explain the factor of two using momentum delivered to the reflector.
\end{enumerate}

\item \textbf{Charge continuity and the outgoing Green function.}
Let $\Jimp$ and $\rho$ be impressed sources on a bounded domain.
\begin{enumerate}[label=(\alph*)]
\item Take the divergence of the harmonic Amp\`ere--Maxwell law and derive Eq.~\eqref{eq:ch1.charge-continuity-harmonic}. Explain why independent specification of $\rho$ and $\Jimp$ is inconsistent.
\item Integrate the continuity equation over a source region $\Omega_s$ and interpret Eq.~\eqref{eq:ch1.global-charge-compatibility} as a global compatibility condition on the impressed current.
\item Using Eqs.~\eqref{eq:ch1.inhomogeneous-helmholtz} and~\eqref{eq:ch1.scalar-green-outgoing}, write the scalar potential at an exterior point $\mathbf{r}$ due to $\rho(\mathbf{r}')$ and identify the outgoing factor $e^{-\jj k_0|\mathbf{r}-\mathbf{r}'|}/(4\pi|\mathbf{r}-\mathbf{r}'|)$. In the far zone, expand $|\mathbf{r}-\mathbf{r}'|$ to first order in $|\mathbf{r}'|/r$ and extract the angular phase factor of Eq.~\eqref{eq:ch1.green-far-zone-phase}.
\item Explain in one paragraph why a Helmholtz decomposition of $\Jimp$ into longitudinal and transverse parts is \emph{not} identical to a reactive-versus-radiative partition of the radiated field \cite{stratton1941,jackson1999}.
\end{enumerate}

\item \textbf{Boundaries, outgoing waves, and uniqueness.}
\begin{enumerate}[label=(\alph*)]
\item Write the PEC, PMC, and scalar passive impedance conditions of Eq.~\eqref{eq:ch1.outer-boundary-classes} in words, including the direction of the outward normal on an exterior problem.
\item For an unbounded radiating exterior, state the outgoing radiation condition of Eq.~\eqref{eq:ch1.radiation-trace-condition} and explain why superposition with an incoming spherical-wave completion is excluded \cite{sommerfeld1949}.
\item \textbf{Lossless media.} For \emph{real} lossless $\varepsilon$ and $\mu$, start from Eq.~\eqref{eq:ch1.uniqueness-local-div}, form the real part of the complex Poynting identity of Eq.~\eqref{eq:ch1.uniqueness-energy-identity}, and integrate over a truncated exterior $\Omega_{\mathrm{ext},R}$ to obtain the energy balance used in the uniqueness argument of Subsection~\ref{sec:ch01.1.3}.
\item \textbf{Passive loss.} For complex passive $\varepsilon=\varepsilon'-\jj\varepsilon''$, $\mu=\mu'-\jj\mu''$ with $\varepsilon'',\mu''\ge0$, use Eq.~\eqref{eq:ch1.uniqueness-local-div-complex}, take its \emph{real part}, and integrate that lossy identity over the domain. Do \emph{not} substitute Eq.~\eqref{eq:ch1.uniqueness-energy-identity} in lossless form when dissipation is present. Summarize what uniqueness follows on: (i)~a bounded dissipative domain with homogeneous PEC/PMC data; (ii)~a lossless cavity at an eigenfrequency; and (iii)~an exterior radiating problem with outgoing selection. State the extra boundary flux for passive impedance surfaces using Eqs.~\eqref{eq:ch1.outer-boundary-classes} and~\eqref{eq:ch1.impedance-boundary-power} \cite{stratton1941,harrington2001}.
\item Explain why compatible impressed sources, interface data, and boundary classes must be fixed before Poynting and torque surface integrals have physical meaning.
\end{enumerate}

\item \textbf{Angular phase from a localized source.}
Using Eq.~\eqref{eq:ch1.green-far-zone-phase} only, show that a compact source distribution imprints a direction-dependent phase $-\jj k_0\hat{\mathbf{r}}\cdot\mathbf{r}'$ on the outgoing field. Compute $\partial(\arg E)/\partial\phi$ for a two-point source model $\Jimp(\mathbf{r}')\propto\delta^3(\mathbf{r}'-\mathbf{r}_1)+\delta^3(\mathbf{r}'-\mathbf{r}_2)$ with $\mathbf{r}_2-\mathbf{r}_1$ transverse to $\hat{\mathbf{r}}$ in the far zone, and explain why this phase gradient is \emph{not} by itself a complete statement of angular-momentum transport \cite{stratton1941}.

\item \textbf{Outgoing cylindrical waves.}
For time-harmonic fields with $\exp(+\jj\omega t)$, write the $\rho$-dependence of $H_m^{(1)}(\kappa\rho)$ and $H_m^{(2)}(\kappa\rho)$ in the limit $\kappa\rho\to\infty$. Identify which Hankel family is outgoing by demanding $e^{+\jj\omega t}$ time dependence and physically outgoing phase. For a single azimuthal mode $e^{\jj m\phi}$, compute the period of $\phi$ and state whether a nonzero $m$ implies a net azimuthal power flux without computing the full Maxwell stress integral \cite{stratton1941,harrington2001}.

\item \textbf{Tangential fields on a sphere.}
Let $\mathbf{f}(\theta,\phi)$ be a tangential vector field on $S^2$.
\begin{enumerate}[label=(\alph*)]
\item Decompose $\mathbf{f}=\nabla_S\psi+\hat{\mathbf{r}}\times\nabla_S\chi$ into surface-gradient and surface-curl contributions (state smoothness assumptions).
\item Explain why a single scalar spherical harmonic $Y_\ell^m(\theta,\phi)$ cannot represent a generic tangential electric field on $S^2$.
\item For a $\hat{\boldsymbol\theta}$-polarized tangential component proportional to $Y_1^0(\theta)$, verify tangentiality and state what additional $\hat{\boldsymbol\phi}$ content is required for a divergence-free tangential field on the sphere \cite{stratton1941,hansen1988}.
\end{enumerate}

\item \textbf{Rotation, helicity, and modal limits.}
\begin{enumerate}[label=(\alph*)]
\item Give one physical example of \emph{active} source rotation and one of \emph{passive} coordinate relabeling (Section~\ref{sec:ch01.3}) and state which observables distinguish them.
\item Under the separable assumptions of Remark~\ref{rem:ch1.m-plus-sigma}, starting from Eq.~\eqref{eq:ch1.ex61-flux-from-stress}, derive the modal bookkeeping $(m+\sigma)P/\omega$. Then state explicitly why this relation is not a universal Maxwell-stress theorem for arbitrary full-vector fields, citing the uniform plane-wave disk countercheck of Subsection~\ref{sec:ch01.6.1}.
\item Using Eq.~\eqref{eq:ch1.helicity-density}, verify that $\int_V \mathbf{e}\cdot\mathbf{h}\,dV$ built from $(\E,\H)$ is unchanged under $\mathbf{A}\to\mathbf{A}+\nabla\chi$, $\phi\to\phi-\jj\omega\chi$, while a potential-based canonical spin density may change \cite{belinfante1940,jackson1999}.
\item Two equal-power outgoing modes with azimuthal indices $+m$ and $-m$ on an aperture produce a standing azimuthal intensity pattern. Using Eq.~\eqref{eq:ch1.ex62-net-am-transport} under the stated modal assumptions, compute the net signed orbital flux and explain why the standing pattern does not reveal that flux by inspection alone \cite{allen1992,beth1936}.
\end{enumerate}

\item \textbf{Fraunhofer condition and aperture phase.}
Starting from the far-zone Green-function phase of Eq.~\eqref{eq:ch1.green-far-zone-phase}, show that for a source confined to a sphere of radius $a$ the Fraunhofer condition $r\gg a^2/\lambda$ makes the curvature term $|\mathbf{r}'|^2/r$ negligible in the exponent. For a uniform circular aperture of radius $a$ in the $z=0$ plane, write the far-field $\mathbf{E}(\mathbf{r})$ as an integral over $\rho'\le a$ with integrand $e^{-\jj k_0 r}e^{-\jj k_0(\rho'^2/2r)}e^{-\jj k_0(\rho'\rho\sin\gamma)/r}$ and identify the role of $\hat{\mathbf{r}}\cdot\mathbf{r}'$ in selecting the observation direction \cite{bornwolf1999,silver1949}.

\item \textbf{Hertzian dipole: fields, power, and shells.}
Take $\mathbf{p}=p_0\hat{\mathbf{z}}$ with $|p_0|$ in C\,m, $k_0=\omega\sqrt{\varepsilon_0\mu_0}$, and Eqs.~\eqref{eq:ch1.dipole-fields-explicit}--\eqref{eq:ch1.dipole-reactive-ratio}.
\begin{enumerate}[label=(\alph*)]
\item From the retarded vector potential with outgoing Green function, show that the far-zone $\mathbf{E}$ scales as $k_0^2 p_0\sin\theta\,e^{-\jj k_0 r}/r$ without re-deriving the full multipole ladder.
\item Use charge continuity to explain why a localized time-harmonic source cannot sustain a radiating electric monopole while a static charge distribution can.
\item For $r>0$, verify $\nabla\cdot\E=0$, $\curl\E=-\jj\omega\mu_0\H$, and $\curl\H=+\jj\omega\varepsilon_0\E$ in spherical coordinates. Identify which terms scale as $r^{-3}$, $r^{-2}$, and $r^{-1}$.
\item At $f=300\,\mathrm{MHz}$ ($\lambda_0=1\,\mathrm{m}$, $k=2\pi\,\mathrm{rad/m}$), compute $kr_1$ and $kr_2$ for $r_1=0.05\,\mathrm{m}$ and $r_2=5\,\mathrm{m}$. Name the dominant radial field terms on each shell, evaluate $|P_{\mathrm{reac}}(r_i)|/P_{\mathrm{rad}}$ from Eq.~\eqref{eq:ch1.dipole-reactive-ratio}, and explain why $P_{\mathrm{rad}}$ is the same on both shells while the reactive exchange is not. \emph{This is a near-zone example at $kr_1\approx0.314$; it is not the asymptotic limit $kr\to0$.}
\item Show $D(\theta)=(3/2)\sin^2\theta$ for the aligned dipole, verify $\int_{4\pi}D(\theta,\phi)\,d\Omega=4\pi$, and state $D_{\max}$.
\end{enumerate}

\item \textbf{Reactive exchange, Chu $Q$, and finite size.}
\begin{enumerate}[label=(\alph*)]
\item Starting from Eq.~\eqref{eq:ch1.ex64-power-functionals}, prove $P_{\mathrm{reac}}(r)=\Im\{\oint_{S_r}\mathbf{S}_\mathrm{c}\cdot\hat{\mathbf{r}}\,dS\}$ and explain why its magnitude depends on $r$ while $P_{\mathrm{rad}}$ does not in a lossless exterior.
\item Using Eq.~\eqref{eq:ch1.chu-q-adopted}, distinguish $Q_{\mathrm{vol}}=2\omega\max(W_e,W_m)/P_{\mathrm{rad}}$ from $P_{\mathrm{reac}}(r)/P_{\mathrm{rad}}$ and from Chu's modal $Q_1$ of Eq.~\eqref{eq:ch1.chu-Qn-def}. Give one sentence on why total field energy in a finite volume generally includes propagating power \cite{chu1948,mcclean1996}.
\item For the dipole example of Problem~9(d), compare reactive dominance at $kr_1=0.1\pi$ with radiative dominance at $kr_2=10\pi$ using Eq.~\eqref{eq:ch1.dipole-reactive-ratio}.
\item From Eq.~\eqref{eq:ch1.chu-q-lower-bound}, state the \emph{exact} ideal $TM_{01}$ radiation $Q$ in McLean's spherical model and expand it for $ka\ll1$ through orders $(ka)^{-3}$ and $(ka)^{-1}$. State explicitly that it is a minimum-$Q$ benchmark for the stated antenna class, not a universal identity for arbitrary radiators \cite{mcclean1996,chu1948}.
\end{enumerate}
\end{enumerate}
